\section*{Chapitre \ref{ch:nw:switches-and-layer-1-devices} --- Commutateurs et dispositifs de couche 1}

\begin{solution}{exo:nw:switches-and-layer-1-devices:1}
Store-and-forward~: $8 \times 512 / 25 + 300 = 163.8 + 300 = \qty{463.8}{\nano\second}$. Cut-through~:
$8 \times 64 / 25 + 300 = 20.5 + 300 = \qty{320.5}{\nano\second}$.
\end{solution}

\begin{solution}{exo:nw:switches-and-layer-1-devices:2}
Depuis 25 en entrée vers 10 en sortie, il peut passer en cut-through~: les bits arrivent plus vite qu'ils ne repartent, de sorte que la sortie ne se vide jamais (la différence
s'accumule dans le tampon jusqu'à la fin de la trame). Depuis 10 en entrée vers 25 en sortie, ce n'est pas possible~: la sortie enverrait plus vite que les bits
n'arrivent et se viderait au milieu de la trame, de sorte que la trame doit d'abord être stockée.
\end{solution}

\begin{solution}{exo:nw:switches-and-layer-1-devices:3}
$2 \times 6 = \qty{12}{\giga\bit\per\second}$ offerts à un port de \qty{10}{\giga\bit\per\second}~: $1 - 10/12 = 1/6$, soit environ 17~\% de
la copie perdue tant que les deux sens sont aussi chargés.
\end{solution}

\begin{solution}{exo:nw:switches-and-layer-1-devices:4}
$D$ a un écart type de $100\sqrt2 = 141.4$~ns~: $2\Phi(96/141.4) - 1 = 2\Phi(0.679) - 1 = 0.50$.
\end{solution}

\begin{solution}{exo:nw:switches-and-layer-1-devices:5}
Sérialisation $8 \times 104/10 = 83.2$~ns, fibre $36 \times 4.88 = 175.6$~ns, dispositifs 4~ns. La fibre~: des câbles plus courts dans la cage
et, surtout, une interconnexion plus courte, ce qui relève de la plateforme (chapitre~9)~; la sérialisation ensuite, grâce à une liaison plus rapide.
\end{solution}

\begin{solution}{exo:nw:switches-and-layer-1-devices:6}
Un fan-out copie une entrée vers des sorties qui ne transportent rien d'autre~: une sortie n'est jamais occupée au moment où arrivent ses bits, de sorte que rien n'attend. La
sortie d'un multiplexeur est partagée par plusieurs entrées, et une trame qui arrive pendant qu'une autre est en cours d'émission doit attendre. La
latence du fan-out est une constante~; celle du multiplexeur est une constante plus un délai de file d'attente qui dépend des autres entrées, sa gigue.
\end{solution}

\begin{solution}{exo:nw:switches-and-layer-1-devices:7}
Chaque multiplexeur sert maintenant deux serveurs~: l'attente moyenne à \qty{50}{\nano\second} tombe de 93.2~ns (quatre sur un seul) à 22.7~ns (deux
sur chacun), valeur de \texttt{mux\_contention(2, 50)}~; la conception gagne un handoff et un multiplexeur.
\end{solution}

\begin{solution}{exo:nw:switches-and-layer-1-devices:8}
Un seul boîtier a une seule alimentation, un seul châssis, un seul jeu d'optiques et une seule configuration~: une panne d'alimentation ou un changement de configuration touche
les deux lignes à la fois. Placer la ligne A et la ligne B sur des dispositifs différents, comme le prescrit la méthode~; l'absence de logiciel élimine un mode de panne,
mais pas les autres.
\end{solution}

\begin{solution}{pb:nw:switches-and-layer-1-devices:1}
\begin{enumerate}
\item Commodité $1\,439.8 + 1\,430.2 = 2\,870.0$~ns~; cut-through $560.0 + 556.8 = 1\,116.7$~ns~; couche 1 $262.8 + 294.6 = 557.3$~ns.
\item 1\,753.3~ns~; puis 559.4~ns.
\item Deux sérialisations de la trame de 104~octets, $2 \times 83.2 = 166.4$~ns, plus \qty{1}{\micro\second} de temps interne du modèle.
\item $36 \times 4.88 = 175.6$~ns.
\item 93.2~ns.
\item 186.1~ns.
\item Oui~: $294.6 + 186.1 = 480.7$~ns, contre 556.8~ns pour le chemin d'ordres en cut-through.
\item Répartis sur les deux multiplexeurs, deux chacun, avec les deux stratégies les plus corrélées sur des multiplexeurs différents.
\item Les deux commutateurs, le câble qui les relie et les deux câbles vers les serveurs.
\item Chaque ligne a son propre commutateur et chaque commutateur atteint les deux serveurs et son propre handoff d'ordres~: la perte d'un commutateur ne fait perdre qu'une ligne
  et un chemin d'ordres, jamais la totalité de l'un ou de l'autre.
\item Deux handoffs d'ordres, chacun sur son propre multiplexeur, et la plateforme acceptant les ordres sur les deux (deux sessions).
\item Rien pour les serveurs~: ils prennent l'autre ligne, qui a la même latence~; la société perd sa protection contre la panne suivante jusqu'au remplacement du dispositif.
  \\
\item $60\,000 / 1\,753.3 \approx 34.2$ par nanoseconde et par an.
\item $90\,000 / 559.4 \approx 160.9$ par nanoseconde et par an.
\item \textbf{Résultat retenu.} Chemins de déclenchement de 2\,870, 1\,117 et 557~ns~; les premières 1\,753 nanosecondes coûtent environ 34 par an chacune, les
  559 suivantes environ 161 chacune.
\item Un système qui gagne ou perd des courses décidées en centaines de nanosecondes sur cette plateforme, comme les stratégies d'arbitrage de latence et de market making
  du livre~11.
\item Filtrage par adresse ou par groupe, routage entre segments, files par port et compteurs d'un commutateur~; la conception doit assurer elle-même le filtrage sur le serveur.
  \\
\item Des câbles (la longueur de l'interconnexion), de la sérialisation (liaisons plus rapides) et des serveurs eux-mêmes (chapitres~3
  et~8, et livre~13).
\item La vérification à point unique, la vérification par tap et une capture des chemins réels à comparer aux chemins calculés (chapitre~5).
\item Celui qui ne lit pas du tout la trame.
\end{enumerate}
\end{solution}

\begin{solution}{iq:nw:switches-and-layer-1-devices:1}
Le store-and-forward reçoit la trame entière, la vérifie, puis l'envoie~: il ajoute une sérialisation complète. Le cut-through commence à émettre
après l'en-tête~: il ajoute une constante. Il stocke puis transmet lorsque la sortie est occupée ou plus rapide que l'entrée, et sur les commutateurs
qui doivent réécrire la trame.
\iqlookfor{le terme de sérialisation et les conditions de repli.}
\end{solution}

\begin{solution}{iq:nw:switches-and-layer-1-devices:2}
Un dispositif qui relie des ports au niveau du signal, régénère les bits sans décoder les trames~: des nanosecondes, pas de files, duplication et brassage
pilotés par logiciel. Il ne peut ni lire les adresses, ni filtrer les groupes, ni router, ni fusionner le trafic de plusieurs ports (un
multiplexeur, qui est un autre dispositif, le fait avec une file).
\iqlookfor{pas de découpage en trames, donc pas de filtrage ni de contention.}
\end{solution}

\begin{solution}{iq:nw:switches-and-layer-1-devices:3}
Les ordres vont de plusieurs vers un~; un fan-out va d'un vers plusieurs. Un multiplexeur les fusionne avec une petite file~: lorsque plusieurs serveurs répondent au même
événement dans le temps de sérialisation d'une trame, leurs ordres se font concurrence les uns les autres, et les stratégies les plus rapides d'une société entrent en
collision précisément quand elles comptent. Répartir les stratégies corrélées sur plusieurs handoffs.
\iqlookfor{une contention corrélée aux événements de marché.}
\end{solution}

\begin{solution}{iq:nw:switches-and-layer-1-devices:4}
Un tap~: il est sur la liaison, copie chaque bit, erreurs comprises, en même temps que l'original, ne peut pas perdre la copie, et continue
de fonctionner lorsque le commutateur ne le fait pas. Un port de recopie est une sortie de commutateur qui peut être surabonnée et dont la copie traverse
les files propres au commutateur, de sorte que sa temporisation n'est pas celle du câble.
\iqlookfor{surabonnement et fidélité temporelle de la copie.}
\end{solution}

\begin{solution}{iq:nw:switches-and-layer-1-devices:5}
Ligne A et ligne B passant par des taps distincts vers des fan-outs distincts, chacun desservant les quatre serveurs~; les ordres de chaque serveur passent par un
multiplexeur vers le handoff, et un second handoff et un second multiplexeur si la plateforme le permet~; capture sur chaque tap~; administration hors bande.
Pannes~: un fan-out (les serveurs utilisent l'autre ligne), un multiplexeur (les serveurs utilisent le second handoff, ou s'arrêtent s'il n'y en a pas~: le
point unique à éliminer par négociation), un câble, un serveur (ses stratégies migrent).
\iqlookfor{aucun dispositif partagé entre A et B, et un énoncé honnête de ce qui reste unique.}
\end{solution}

\begin{solution}{iq:nw:switches-and-layer-1-devices:6}
À l'ouverture, les ports de sortie sont occupés~: les trames qui trouvent leur sortie en train d'émettre une autre trame sont stockées et attendent, de sorte que le commutateur
ajoute de la file d'attente (et perd le cut-through pour ces trames). Regarder le taux d'occupation de la \omterm{def:nw:networking-for-trading:egress}{file de sortie} et les \omterm{def:nw:networking-for-trading:burst}{micro-rafales} (chapitre~1), et non
la latence à vide de la fiche technique.
\iqlookfor{file d'attente dépendant de la charge et repli vers le store-and-forward.}
\end{solution}
